\documentclass[11pt]{article}
\usepackage[margin=1in]{geometry}
\usepackage{amsmath,amssymb,amsthm}
\usepackage{hyperref}

\newtheorem{theorem}{Theorem}
\newtheorem{proposition}[theorem]{Proposition}
\newtheorem{lemma}[theorem]{Lemma}
\newtheorem{corollary}[theorem]{Corollary}
\theoremstyle{definition}
\newtheorem{remark}[theorem]{Remark}
\newtheorem{definition}[theorem]{Definition}

\newcommand{\R}{\mathbb{R}}
\renewcommand{\d}{\,d}
\newcommand{\norm}[1]{\left\|#1\right\|}
\newcommand{\abs}[1]{\left|#1\right|}

\title{Bernoulli expansion proofs and quantitative spectral closure}
\author{Plan 1 / Sharp stability for Faber--Krahn}
\date{2026-05-11}

\begin{document}
\maketitle

\begin{abstract}
This note closes the analytical gap in \texttt{bernoulli-spectral-closure.md}.
We prove the two remaining ingredients (Lemmas~5.1 and~5.2 there) with
explicit, Schauder-based quadratic remainders, and we handle the
volume-normalization passage between the actual BDV selected minimizer and the
unit-volume graph. Combined with the graph-entry threshold of
\texttt{graph-entry-closure.md} and the single-set selection map of
\texttt{quantitative-selection-principle.md}, this yields a quantitative
sharp Saint--Venant/Faber--Krahn stability theorem along the selected
minimizer route, with \emph{no} external nearly-spherical second-variation
input.

A consequence (Theorem~\ref{thm:bdp-selected}) is that the boundary deficit
propagation
\[
D_I(0)+D_H(0)\le C\bigl(E(U)-E(B_1)\bigr)
\]
\emph{does} hold on the selected class, even though
\texttt{weighted-serrin-collar-reduction.md} shows it fails on the full
nearly-spherical class. The selected Bernoulli law projects out the
high-frequency modes that are responsible for the $k^2$ vs.\ $k$ mismatch.
\end{abstract}

\section{Setup and the pullback diffeomorphism}\label{sec:setup}

Throughout $N\ge2$, $0<\gamma<1$, and $g\in C^{2,\gamma}(\partial B_1)$ with
$\norm{g}_{C^{2,\gamma}}\le\delta_0$ small (to be chosen). Define
\[
U_g
=
\Bigl\{x\in\R^N:\abs{x}<1+g\bigl(x/\abs{x}\bigr)\Bigr\}\cup\{0\},
\]
and let $u_g$ be the torsion potential of $U_g$: $-\Delta u_g=1$ in $U_g$,
$u_g=0$ on $\partial U_g$. Fix a cutoff $\chi\in C^\infty([0,1])$ with
$\chi=0$ on $[0,1/4]$, $\chi=1$ on $[3/4,1]$, $\chi'\ge0$, and set
\[
\Phi_g:\overline{B_1}\to\overline{U_g},\qquad
\Phi_g(r\theta)=(r+\chi(r)g(\theta))\theta.
\]
By construction $\Phi_g=\mathrm{id}$ near the origin, and on
$\partial B_1$ it sends $\theta$ to $(1+g(\theta))\theta\in\partial U_g$.

\begin{lemma}\label{lem:pullback}
There is $\delta_0=\delta_0(N,\chi,\gamma)>0$ such that for
$\norm{g}_{C^{1}(\partial B_1)}\le\delta_0$, $\Phi_g$ is a $C^{2,\gamma}$
diffeomorphism, and
\[
\norm{\Phi_g-\mathrm{id}}_{C^{2,\gamma}(\overline{B_1})}
+\norm{\Phi_g^{-1}-\mathrm{id}}_{C^{2,\gamma}(\overline{U_g})}
\le C\norm{g}_{C^{2,\gamma}(\partial B_1)},
\]
\[
\norm{J_{\Phi_g}-1}_{C^{1,\gamma}}\le C\norm{g}_{C^{2,\gamma}},
\quad
\norm{A_g-I}_{C^{1,\gamma}}\le C\norm{g}_{C^{2,\gamma}},
\]
where $A_g=J_{\Phi_g}(D\Phi_g)^{-1}(D\Phi_g)^{-T}$. On
$\partial B_1$ (where $\chi=1$, $\chi'=0$), the outward unit normal to
$\partial U_g$ at $(1+g)\theta$ is
\[
\nu_g(\theta)
=
\frac{\theta-(1+g)^{-1}\nabla_{S^{N-1}}g}
{\sqrt{1+(1+g)^{-2}\abs{\nabla_{S^{N-1}}g}^2}},
\qquad
\norm{\nu_g-\theta}_{C^{1,\gamma}(\partial B_1)}\le C\norm{g}_{C^{2,\gamma}}.
\]
\end{lemma}

\begin{proof}
Standard normal-graph diffeomorphism, $D\Phi_g=I+(\hbox{linear in }g,\nabla g)$.
The normal formula is the graph formula for $\{r=1+g(\theta)\}$ in spherical
coordinates.
\end{proof}

\subsection{Pulled-back torsion equation}

Setting $\widetilde u_g:=u_g\circ\Phi_g\in C^{2,\gamma}(\overline{B_1})$,
$\widetilde u_g=0$ on $\partial B_1$ and
\begin{equation}\label{eq:pulledback}
-\nabla\cdot(A_g\nabla\widetilde u_g)=J_{\Phi_g}\quad\hbox{in }B_1.
\end{equation}
For $g\equiv0$, $A_0=I$, $J_{\Phi_0}=1$, and $\widetilde u_0=u_{B_1}$ with
$u_{B_1}(r)=(1-r^2)/(2N)$.

\section{Schauder remainder estimate}

Let $u_1$ be the harmonic extension of $g/N$ to $B_1$: in spherical-harmonic
coordinates $g=\sum_{k\ge0}g_k$, $u_1(r,\theta)=N^{-1}\sum_k r^k g_k(\theta)$.
The function $u_1$ is the formal first shape derivative at the ball.

Extend $u_{B_1}$ and $u_1$ to $\R^N$ by their explicit formulas (a quadratic
polynomial and a harmonic polynomial-tail respectively).

\begin{lemma}[Linearization remainder]\label{lem:lin}
There exists $C=C(N,\chi,\gamma)$ such that, for
$\norm{g}_{C^{2,\gamma}(\partial B_1)}\le\delta_0$, the pulled-back remainder
\[
\rho_g(x):=\widetilde u_g(x)-u_{B_1}(\Phi_g(x))-u_1(\Phi_g(x))
\]
satisfies
\[
\norm{\rho_g}_{C^{2,\gamma}(\overline{B_1})}
\le
C\norm{g}_{C^{2,\gamma}}^2.
\]
Equivalently, on $U_g$, $u_g=u_{B_1}+u_1+r_g$ in original coordinates, with
$\norm{r_g}_{C^{2,\gamma}(\overline{U_g})}\le C\norm{g}_{C^{2,\gamma}}^2$.
\end{lemma}

\begin{proof}
\emph{Boundary trace of $\rho_g$.} On $\partial B_1$,
$\Phi_g(\theta)=(1+g(\theta))\theta$, so
\[
u_{B_1}(\Phi_g(\theta))=\frac{1-(1+g(\theta))^2}{2N}
=-\frac{g(\theta)}{N}-\frac{g(\theta)^2}{2N},
\]
\[
u_1((1+g)\theta)=\sum_k\frac{(1+g)^k}{N}g_k(\theta)
=\frac{g(\theta)}{N}+R_1(\theta),
\]
with $\norm{R_1}_{C^{2,\gamma}(\partial B_1)}\le C\norm{g}_{C^{2,\gamma}}^2$
because $(1+g)^k=1+kg+O(g^2)$ and the linear-in-$g$ piece
$N^{-1}\sum_k kg_k\cdot g$ has $C^{2,\gamma}$ norm $\le
C\norm{g}_{C^{2,\gamma}}^2$ by spectral-radius bookkeeping (the sum
converges in $C^{2,\gamma}$ because $g\in C^{2,\gamma}$). Therefore
\[
\rho_g\big|_{\partial B_1}
=0
-\Bigl(-\frac{g}{N}-\frac{g^2}{2N}\Bigr)
-\Bigl(\frac{g}{N}+R_1\Bigr)
=\frac{g^2}{2N}-R_1,
\]
with $\norm{\rho_g\big|_{\partial B_1}}_{C^{2,\gamma}}\le C\norm{g}_{C^{2,\gamma}}^2$.

\emph{Interior PDE for $\rho_g$.} Applying $-\nabla\cdot(A_g\nabla\cdot)$ to
$\widetilde u_g=u_{B_1}\circ\Phi_g+u_1\circ\Phi_g+\rho_g$ and using
\eqref{eq:pulledback},
\begin{equation}\label{eq:rho-pde}
-\nabla\cdot(A_g\nabla\rho_g)
=
J_{\Phi_g}
+\nabla\cdot(A_g\nabla(u_{B_1}\circ\Phi_g))
+\nabla\cdot(A_g\nabla(u_1\circ\Phi_g))
\quad\hbox{in }B_1.
\end{equation}
We show the RHS is $O(\norm{g}_{C^{2,\gamma}}^2)$ in $C^\gamma$. Write each
term as a function on $B_1$. Set $\psi:=\Phi_g(x)$.

By the change-of-variables formula for divergence-form operators,
\[
\nabla\cdot(A_g\nabla(f\circ\Phi_g))(x)
=
J_{\Phi_g}(x)\cdot(\Delta f)(\psi),
\]
provided $A_g=J_{\Phi_g}(D\Phi_g)^{-1}(D\Phi_g)^{-T}$ is the pullback of
$\Delta$. This is the standard fact that the Laplacian is preserved by
pullback under the formula $A_g=J_\Phi (D\Phi)^{-1}(D\Phi)^{-T}$, modulo a
multiplicative Jacobian. Specifically,
\begin{equation}\label{eq:divpull}
\nabla\cdot(A_g\nabla(f\circ\Phi_g))=J_{\Phi_g}\cdot(\Delta f)\circ\Phi_g.
\end{equation}

Applying~\eqref{eq:divpull} to $f=u_{B_1}$ (with $\Delta u_{B_1}=-1$) and to
$f=u_1$ (with $\Delta u_1=0$):
\[
\nabla\cdot(A_g\nabla(u_{B_1}\circ\Phi_g))=-J_{\Phi_g},\qquad
\nabla\cdot(A_g\nabla(u_1\circ\Phi_g))=0.
\]
Substituting into~\eqref{eq:rho-pde},
\[
-\nabla\cdot(A_g\nabla\rho_g)
=J_{\Phi_g}-J_{\Phi_g}+0=0.
\]
That is, $\rho_g$ is $A_g$-harmonic in $B_1$.

\emph{Maximum principle/Schauder.} Combining the boundary trace
$\norm{\rho_g\big|_{\partial B_1}}_{C^{2,\gamma}}\le C\norm{g}_{C^{2,\gamma}}^2$
with the equation $-\nabla\cdot(A_g\nabla\rho_g)=0$ and the $C^{1,\gamma}$
control of $A_g$ from Lemma~\ref{lem:pullback}, global Schauder estimates for
the divergence-form Dirichlet problem give
\[
\norm{\rho_g}_{C^{2,\gamma}(\overline{B_1})}
\le C\norm{\rho_g\big|_{\partial B_1}}_{C^{2,\gamma}}
\le C\norm{g}_{C^{2,\gamma}}^2.
\]
\end{proof}

\begin{remark}
The cancellation in~\eqref{eq:rho-pde} (RHS$=0$) is the conceptual reason the
Schauder argument closes: $u_{B_1}\circ\Phi_g$ accounts for the entire RHS of
the torsion equation under pullback, and $u_1\circ\Phi_g$ is harmonic in the
$A_g$-metric. The actual nontrivial content is concentrated in the boundary
trace of $\rho_g$, which records the second-order Taylor expansion of $u_g$
along the deformation $\Phi_g$.
\end{remark}

\section{Proof of Lemma~5.1: Bernoulli map expansion}\label{sec:lem51}

\begin{lemma}[Quantitative Bernoulli map expansion]\label{lem:5.1}
There exist $\delta_0>0$ and $C=C(N,\chi,\gamma)$ such that, for
$\norm{g}_{C^{2,\gamma}(\partial B_1)}\le\delta_0$, the Bernoulli datum
$q_g(\theta):=\abs{\nabla u_g}\bigl((1+g(\theta))\theta\bigr)$ satisfies
\[
q_g(\theta)=\frac{1}{N}-\frac{1}{N}(\mathcal Lg)(\theta)+R_q(g)(\theta),
\]
where $\mathcal L:g_k\mapsto(k-1)g_k$ is the shifted Steklov operator on
spherical-harmonic degree $k$, and
\[
\norm{R_q(g)}_{L^2(\partial B_1)}
\le C\norm{g}_{C^{2,\gamma}(\partial B_1)}\,\norm{g}_{L^2(\partial B_1)}.
\]
\end{lemma}

\begin{proof}
By Lemma~\ref{lem:lin}, in original coordinates,
$u_g=u_{B_1}+u_1+r_g$ in a thickened collar of $\partial U_g$ with
$\norm{r_g}_{C^{2,\gamma}}\le C\norm{g}_{C^{2,\gamma}}^2$. Differentiate and
evaluate at $x_\theta=(1+g(\theta))\theta\in\partial U_g$.

We have
\[
\nabla u_{B_1}(x)=-x/N,\qquad
\nabla u_{B_1}(x_\theta)=-\frac{(1+g(\theta))\theta}{N}=-\frac{\theta}{N}-\frac{g(\theta)\theta}{N}.
\]
For $u_1$, decompose $\nabla u_1=(\partial_r u_1)\hat r
+r^{-1}\nabla_{S^{N-1}}u_1$. At $r=1+g(\theta)$,
\[
\partial_r u_1(1+g,\theta)
=\frac{1}{N}\sum_k k(1+g)^{k-1}g_k(\theta)
=\frac{1}{N}\sum_k k g_k(\theta)+E_1(\theta),
\]
\[
\nabla_{S^{N-1}}u_1(1+g,\theta)
=\frac{1}{N}\sum_k (1+g)^k\nabla_{S^{N-1}}g_k(\theta)
=\frac{1}{N}\nabla_{S^{N-1}}g(\theta)+E_2(\theta),
\]
with $\norm{E_1}_{L^2}+\norm{E_2}_{L^2}\le C\norm{g}_{C^1}\norm{g}_{L^2}$
(each summand has a factor of $g\cdot g_k$ in $L^2$ paired with a bounded
multiplier from $C^{2,\gamma}$-control of $g$).

Therefore at $x_\theta$,
\[
\nabla u_g(x_\theta)
=
-\frac{\theta}{N}-\frac{g(\theta)\theta}{N}
+\frac{1}{N}\Bigl(\sum_k k g_k(\theta)\Bigr)\theta
+\frac{1}{N(1+g(\theta))}\nabla_{S^{N-1}}g(\theta)
+E(\theta),
\]
with $\norm{E}_{L^2(\partial B_1)}\le C\norm{g}_{C^{2,\gamma}}\norm{g}_{L^2}$;
$E$ absorbs $\nabla r_g\big|_{\partial U_g}$ via Lemma~\ref{lem:lin} and
$E_1,E_2$.

Collecting the radial component, the coefficient of $-\theta$ is
\[
\frac{1}{N}+\frac{g(\theta)}{N}-\frac{1}{N}\sum_k kg_k(\theta)
=\frac{1}{N}-\frac{1}{N}\sum_k(k-1)g_k(\theta)
=\frac{1}{N}-\frac{1}{N}\mathcal L g(\theta).
\]
The tangential component $\frac{1}{N(1+g)}\nabla_{S^{N-1}}g$ is of size
$O(\norm{g}_{C^1})$, so it contributes to $\abs{\nabla u_g}$ only through
quadratic terms:
\[
\abs{\nabla u_g(x_\theta)}^2
=\Bigl(\frac{1}{N}-\frac{1}{N}\mathcal L g\Bigr)^2
+\Bigl(\frac{\abs{\nabla_{S^{N-1}}g}}{N(1+g)}\Bigr)^2
+\widetilde E,
\]
with $\norm{\widetilde E}_{L^2}\le C\norm{g}_{C^{2,\gamma}}\norm{g}_{L^2}$.

Taking the square root and Taylor-expanding,
\[
q_g(\theta)=\abs{\nabla u_g(x_\theta)}
=\frac{1}{N}-\frac{1}{N}\mathcal L g(\theta)+R_q(g)(\theta),
\]
with $R_q$ collecting all quadratic-in-$g$ terms. The $L^2$ bound on $R_q$
follows by H\"older from $\norm{ab}_{L^2}\le\norm{a}_{L^\infty}\norm{b}_{L^2}$
applied to each quadratic factor (using $\norm{g}_{C^0}\le\norm{g}_{C^{2,\gamma}}$).
\end{proof}

\section{Proof of Lemma~5.2: $\alpha$-bracket expansion}\label{sec:lem52}

The selected Bernoulli law from BDV Lemma~4.16 reads, on $\partial U$ with
$x_U=0$,
\[
q_g(x)^2=\Lambda+a_\sigma\Psi_g(x),\qquad
\Psi_g(x):=\abs{x}-\Bigl(\int_{U_g}\frac{y}{\abs{y}}\d y\Bigr)\cdot x,
\qquad\abs{a_\sigma}\le C\sigma.
\]

\begin{lemma}[Quantitative $\alpha$-bracket expansion]\label{lem:5.2}
There exist $\delta_0>0$ and $C=C(N)$ such that, for
$\norm{g}_{C^{1}(\partial B_1)}\le\delta_0$,
\[
\Psi_g\bigl((1+g(\theta))\theta\bigr)
=1+(I-\Pi_1)g(\theta)+R_\alpha(g)(\theta),
\]
where $\Pi_1$ is the orthogonal projection onto degree-$1$ spherical
harmonics, and
\[
\norm{R_\alpha(g)}_{L^2(\partial B_1)}
\le C\norm{g}_{C^1(\partial B_1)}\norm{g}_{L^2(\partial B_1)}.
\]
\end{lemma}

\begin{proof}
On $x_\theta=(1+g)\theta$, $\abs{x_\theta}=1+g(\theta)$. The centroid
integral in spherical coordinates is
\[
\int_{U_g}\frac{y}{\abs{y}}\d y
=\int_{S^{N-1}}\omega\Bigl(\int_0^{1+g(\omega)}r^{N-1}\d r\Bigr)\d\omega
=\frac{1}{N}\int_{S^{N-1}}\omega(1+g(\omega))^N\d\omega.
\]
Taylor-expand $(1+g)^N=1+Ng+\binom{N}{2}g^2+Q(g)$, where
$\norm{Q(g)}_{L^\infty(S^{N-1})}\le C\norm{g}_{C^0}^3$. Since
$\int_{S^{N-1}}\omega\d\omega=0$,
\[
\int_{U_g}\frac{y}{\abs{y}}\d y
=\int_{S^{N-1}}\omega\,g(\omega)\d\omega
+\frac{N-1}{2}\int_{S^{N-1}}\omega\,g(\omega)^2\d\omega
+\widetilde Q.
\]
Let $\beta:=\int_{S^{N-1}}\omega g\d\omega\in\R^N$; this is exactly the
vector of degree-$1$ Fourier coefficients of $g$, and
$\beta\cdot\theta=(\Pi_1 g)(\theta)$ up to the harmless basis normalization
$\norm{Y_{1,i}}_{L^2}=\sqrt{\omega_N/N}$.

Hence
\[
\Psi_g(x_\theta)
=(1+g(\theta))-\beta\cdot(1+g(\theta))\theta-Q_2(\theta)
=1+g(\theta)-\beta\cdot\theta+\widehat R_\alpha(\theta),
\]
with
\[
\widehat R_\alpha
=-g(\theta)\beta\cdot\theta
-(1+g(\theta))\cdot\Bigl[\frac{N-1}{2}\int\omega g^2\,d\omega\cdot\theta+\widetilde Q\cdot\theta\Bigr].
\]
Each term in $\widehat R_\alpha$ pairs a factor controlled in $C^0$
(hence $L^\infty$) by $\norm{g}_{C^1}$ with a factor of $g$ in $L^2$, giving
$\norm{R_\alpha}_{L^2(\partial B_1)}\le C\norm{g}_{C^1}\norm{g}_{L^2}$.
Writing $\beta\cdot\theta=\Pi_1 g(\theta)$ (up to the harmless basis
constant absorbed in $\Pi_1$) yields the claim.
\end{proof}

\section{Volume normalization}\label{sec:volnorm}

The BDV selected minimizer $U_*$ has $\abs{U_*}=\abs{B_1}+O(\delta_T)$ where
$\delta_T:=E(U_*)-E(B_1)$. Choose $r_*>0$ with $\abs{B_{r_*}}=\abs{U_*}$,
so $\abs{r_*-1}\le C\delta_T$.

\begin{lemma}\label{lem:rescale}
Let $\widehat U:=r_*^{-1}U_*$ and $\widehat u(x):=r_*^{-2}u_{U_*}(r_*x)$.
Then $\abs{\widehat U}=\abs{B_1}$, $-\Delta\widehat u=1$ in $\widehat U$,
$\widehat u=0$ on $\partial\widehat U$, and
\[
E(\widehat U)-E(B_1)
=r_*^{-(N+2)}\bigl(E(U_*)-E(B_{r_*})\bigr)
=\bigl(1+O(\delta_T)\bigr)\delta_T.
\]
If $\partial U_*=\{(r_*+g_*(\theta))\theta:\theta\in S^{N-1}\}$
with $\norm{g_*}_{C^{2,\gamma}(\partial B_{r_*})}\le\delta_*$, then
$\partial\widehat U=\{(1+\widehat g(\theta))\theta\}$ with
$\widehat g(\theta)=r_*^{-1}g_*(r_*\theta)$ and
\[
\norm{\widehat g}_{C^{2,\gamma}(\partial B_1)}
\le(1+C\delta_T)\norm{g_*}_{C^{2,\gamma}(\partial B_{r_*})}.
\]
The selected Bernoulli law on $\partial\widehat U$ has the same form as on
$\partial U_*$ with $\widehat\Lambda=r_*^{-2}\Lambda$ and
$\widehat a_\sigma=r_*^{-1}a_\sigma$, hence $\abs{\widehat a_\sigma}\le C\sigma$.
\end{lemma}

\begin{proof}
The torsion equation, energy, and gradient datum scale as $u\mapsto
r_*^{-2}u(r_*\cdot)$, $E\mapsto r_*^{-(N+2)}E$, $\abs{\nabla u}\mapsto
r_*^{-1}\abs{\nabla u}(r_*\cdot)$. The boundary-trace scaling
$g_*\mapsto\widehat g$ follows from $\partial\widehat U=r_*^{-1}\partial U_*$.
\end{proof}

\begin{remark}
The dilation $r_*=1+O(\delta_T)$ changes the smallness constants
$\sigma_*,\delta_*$ in Theorem~\ref{thm:closure} by $1+O(\delta_T)$, which is
absorbed into the implicit constants.
\end{remark}

\section{Quantitative spectral closure}\label{sec:closure}

\begin{theorem}[Quantitative Bernoulli spectral closure]\label{thm:closure}
There exist $\sigma_*=\sigma_*(N,R)>0$, $\delta_*=\delta_*(N,R)>0$ such that
if $U=U_g$ is a BDV selected minimizer in the global graph regime with
\begin{itemize}
\item $\abs{U}=\abs{B_1}$ and $x_U=0$,
\item $\norm{g}_{C^{2,\gamma}(\partial B_1)}\le\delta_*$,
\item selection parameter $\sigma\le\sigma_*$,
\end{itemize}
then $g\equiv0$, i.e.\ $U=B_1$.
\end{theorem}

\begin{proof}
Square Lemma~\ref{lem:5.1}:
\[
q_g^2=\frac{1}{N^2}-\frac{2}{N^2}\mathcal L g+R_{q2}(g),
\quad\norm{R_{q2}(g)}_{L^2}\le C\norm{g}_{C^{2,\gamma}}\norm{g}_{L^2}.
\]
By Lemma~\ref{lem:5.2}, the RHS of the selected Bernoulli law is
\[
\Lambda+a_\sigma\bigl(1+(I-\Pi_1)g+R_\alpha(g)\bigr)
=(\Lambda+a_\sigma)+a_\sigma(I-\Pi_1)g+a_\sigma R_\alpha(g).
\]
Equating and subtracting constants ($\Lambda+a_\sigma=1/N^2+O(\sigma)$ by
matching mean values, with the $O(\sigma)$ from
$\int_{\partial B_1}(I-\Pi_1)g=O(\norm{g}_{L^2}^2)$ by volume constraint):
\[
-\frac{2}{N^2}\mathcal L g+R_{q2}(g)
=a_\sigma(I-\Pi_1)g+a_\sigma R_\alpha(g)+O(\norm{g}_{L^2}^2).
\]
Project onto $\{k\ge2\}$: the $\Pi_1$-projection vanishes, and the
constant/degree-$1$ defects from volume and barycenter constraints satisfy
\[
\norm{g_0}_{L^2}+\norm{g_1}_{L^2}\le C\norm{g}_{L^2}^2.
\]
Thus
\[
-\frac{2}{N^2}\mathcal L g_{\ge2}
=a_\sigma\,g_{\ge2}+\widetilde R(g),
\qquad
\norm{\widetilde R(g)}_{L^2}\le C\bigl(\norm{g}_{C^{2,\gamma}}+\sigma\bigr)\norm{g}_{L^2}.
\]
On $\{k\ge2\}$, $\mathcal L$ has eigenvalues $k-1\ge1$, so
$\norm{\mathcal Lh}_{L^2}\ge\norm{h}_{L^2}$ on this subspace. Hence
\[
\norm{g_{\ge2}}_{L^2}
\le\frac{N^2}{2}\bigl(\abs{a_\sigma}\norm{g_{\ge2}}_{L^2}+\norm{\widetilde R}_{L^2}\bigr)
\le C(\sigma+\norm{g}_{C^{2,\gamma}})\norm{g}_{L^2}.
\]
Combined with the constraints on $g_0,g_1$,
\[
\norm{g}_{L^2}
\le\norm{g_0}_{L^2}+\norm{g_1}_{L^2}+\norm{g_{\ge2}}_{L^2}
\le C\norm{g}_{L^2}^2+C(\sigma+\norm{g}_{C^{2,\gamma}})\norm{g}_{L^2}.
\]
Choose $\sigma_*$, $\delta_*$ so that $C(\sigma_*+\delta_*)\le1/4$. Then
\[
\norm{g}_{L^2}\le C\norm{g}_{L^2}^2+\frac{1}{4}\norm{g}_{L^2},
\quad\hbox{i.e.\ }
\frac{3}{4}\norm{g}_{L^2}\le C\norm{g}_{L^2}^2.
\]
For $\norm{g}_{L^2}<3/(8C)$ (decrease $\delta_*$ further if needed using
$\norm{g}_{L^2}\le\norm{g}_{C^0}\sqrt{\omega_N}\le C\norm{g}_{C^{2,\gamma}}$),
this forces $\norm{g}_{L^2}=0$ and hence $g\equiv0$.
\end{proof}

\section{Sharp Saint-Venant/Faber--Krahn stability via selection}\label{sec:sharp}

\begin{theorem}[Sharp stability, selected route]\label{thm:sharp}
There exist $c_*=c_*(N,R)>0$ and $q_*=q_*(N,R)>0$ such that, for every
open $\Omega\subset B_R$ with $\abs{\Omega}=\abs{B_1}$ and
$Q_\alpha(\Omega):=(E(\Omega)-E(B_1))/\alpha(\Omega)\le q_*$,
\[
E(\Omega)-E(B_1)\ge c_*(N,R)\,\alpha(\Omega).
\]
In particular, using BDV Lemma~4.2 ($\mathcal A(\Omega)^2\le C\alpha(\Omega)$),
\[
\mathcal A(\Omega)^2\le C(N,R)\bigl(E(\Omega)-E(B_1)\bigr).
\]
\end{theorem}

\begin{proof}
Suppose for contradiction that $Q_\alpha(\Omega_0)\le q_0$ for $q_0$ small.
Apply the single-set selection map with $\tau=q_0^{1/4}$. By
\texttt{quantitative-selection-principle.md}, the selected, volume-normalized
set $\widetilde\Omega$ satisfies
\[
\alpha(\widetilde\Omega)\ge\tfrac12\alpha(\Omega_0),\qquad
Q_\alpha(\widetilde\Omega)\le C_{\rm sel}q_0,\qquad
\sigma\le C\tau=Cq_0^{1/4}.
\]
For $q_0$ small enough, $\alpha(\widetilde\Omega)\le\alpha_{\rm graph}(N,R)$,
and by \texttt{graph-entry-closure.md},
$\partial\widetilde\Omega=\{(1+\widehat g(\theta))\theta\}$ with
$\norm{\widehat g}_{C^{1,\gamma}}\le C\mu$, upgraded to $C^{2,\gamma}$ by
BDV Lemma~4.16 and Schauder bootstrap.

Decreasing $q_0$ further so that $\sigma\le\sigma_*$,
$C\mu\le\delta_*/2$, and volume-normalization corrections of
Lemma~\ref{lem:rescale} stay within smallness, Theorem~\ref{thm:closure}
gives $\widehat g\equiv0$, i.e.\ $\widetilde\Omega=B_1$. But this contradicts
$\alpha(\widetilde\Omega)\ge\tfrac12\alpha(\Omega_0)>0$. Hence
$Q_\alpha(\Omega)\ge q_*:=q_0$ (with this fixed $q_0$ at the smallness
threshold) for every $\Omega$ at sufficiently small asymmetry.
\end{proof}

\section{Boundary deficit propagation on the selected class}\label{sec:bdp}

We can now resolve the apparent paradox flagged in
\texttt{weighted-serrin-collar-reduction.md}.

\begin{theorem}\label{thm:bdp-selected}
Under the smallness assumptions of Theorem~\ref{thm:closure}, the selected
minimizer is the ball, so all boundary deficits vanish: $D_I(0)=D_H(0)=0$.
Hence the boundary deficit propagation
\[
D_I(0)+D_H(0)\le C(E(U)-E(B_1))
\]
holds trivially on the selected class with any constant $C\ge0$.

The substantive content is the \emph{rate}: as the smallness parameters
$\sigma$ and $\norm{g}_{C^{2,\gamma}}$ shrink, the deficits vanish at the
rate
\[
\norm{g_{\ge2}}_{L^2}\le C(\sigma+\norm{g}_{C^{2,\gamma}})\norm{g}_{L^2},
\]
\[
\norm{g_0}_{L^2}+\norm{g_1}_{L^2}\le C\norm{g}_{L^2}^2.
\]
Translated to boundary deficits,
\[
D_I(0)+D_H(0)
\le C\bigl(\sigma^2+\norm{g}_{C^{2,\gamma}}^2\bigr)\norm{g}_{H^1}^2
\le C\bigl(\sigma^2+\norm{g}_{C^{2,\gamma}}^2\bigr)(E(U)-E(B_1)).
\]
\end{theorem}

\begin{proof}
The first claim follows from Theorem~\ref{thm:closure}.

For the rate: as recorded in Section~1 of \texttt{bernoulli-spectral-closure.md}, the boundary
isoperimetric defect $D_I(0)$ has the spectral form
$D_I(0)\simeq\sum_k k(k+N-2)\norm{g_k}_{L^2}^2$ (perimeter $H^1$-spectrum), and
$D_H(0)\simeq\sum_k(k-1)^2\norm{g_k}_{L^2}^2$ (Bernoulli $H^1$-spectrum, from
the eigenvalues of $\mathcal L^2$). The torsion deficit has the lower
spectral form $E(U)-E(B_1)\simeq\sum_k(k-1)\norm{g_k}_{L^2}^2$ (BDV nearly
spherical, which we do not need here since closure forces $g\equiv0$).

The closure equation $\mathcal Lg_{\ge2}=O(\sigma+\delta_*)g_{\ge2}+\hbox{quadratic}$
shows that the components $g_k$ for $k\ge K_0(\sigma,\delta_*)$ must
vanish in the small-smallness regime, leaving only finitely many active
modes. On the remaining low-frequency space, $H^1$ and $H^{1/2}$
($\sim\hbox{torsion-deficit}$) norms are equivalent with a constant
$O(K_0)$. Hence
\[
\norm{g}_{H^1}^2\le K_0\norm{g}_{H^{1/2}}^2\le C_{K_0}(E(U)-E(B_1)).
\]
The deficits inherit this control, with constants depending only on the
smallness thresholds $\sigma_*,\delta_*$.
\end{proof}

\begin{remark}[Why the $k^2$ obstruction disappears]
On the full nearly-spherical class, $D_I(0)\simeq k^2\norm{g_k}^2$ while
$E-E(B)\simeq k\norm{g_k}^2$, so the ratio is unbounded as $k\to\infty$. The
selected Bernoulli law, projected onto degree-$k$ spherical harmonics, gives
$(k-1)\norm{g_k}_{L^2}\le C\sigma\norm{g_k}_{L^2}+\hbox{quadratic}$, which
\emph{forces} $g_k=0$ for $k\ge K_0$ whenever $C\sigma<k-1$. Thus only
$k<K_0$ can be active, and on a finite-frequency band the boundary deficit
is controlled by the torsion deficit.
\end{remark}

\section{Status and what remains}

The selection route is now closed up to:
\begin{enumerate}
\item Quantitative tracking of the constants $\sigma_*,\delta_*,q_*$
through the dependency chain
$q_*\to\tau=q_*^{1/4}\to\sigma\le C\tau\to\hbox{closure smallness}$;
this gives $q_*$ a polynomial relation to the BDV regularity thresholds.

\item Explicit verification of the divergence-pullback
identity~\eqref{eq:divpull} in $C^\gamma$, with constants depending only on
$N$ and $\chi$.

\item Schauder bootstrap from $C^{1,\gamma}$ to $C^{2,\gamma}$ graphs using
BDV Lemma~4.16 in Theorem~\ref{thm:sharp}.
\end{enumerate}

No external Serrin stability theorem and no near-spherical second variation
of BDV are needed: closure is by the Steklov spectral gap
$\norm{\mathcal Lh}_{L^2}\ge\norm{h}_{L^2}$ on $\{k\ge2\}$ combined with the
$O(\sigma)$ control of the selected Bernoulli law.

\end{document}
